\section{Justificación}

La implementación de una aplicación web integrada al Sistema Integral de Producción representa una respuesta directa y fundamentada a las limitaciones operativas identificadas en el área de producción y ensamblaje de TenarisTamsa.

En primer lugar, la digitalización de los procesos de captura de datos en planta elimina la dependencia de registros manuales en papel, reduciendo significativamente el riesgo de errores humanos, la pérdida de información y los retrasos en la actualización de datos. Al garantizar la disponibilidad de la información en tiempo real, el personal operativo podrá consultar el estado actual de los procesos productivos de manera inmediata y confiable, mejorando la capacidad de respuesta ante contingencias y facilitando la toma de decisiones oportuna.

En segundo lugar, la centralización de la información proveniente de distintas fuentes en una única plataforma resuelve la fragmentación actual de los sistemas, eliminando la necesidad de consultar múltiples fuentes de manera independiente. Esto permite consolidar y analizar la información de forma integrada, brindando trazabilidad completa sobre las órdenes de producción a lo largo de todas las etapas del proceso, lo cual resulta fundamental para la mejora continua y el control operativo.

Asimismo, al desarrollar el sistema bajo un stack tecnológico moderno y estandarizado, se sientan las bases para garantizar la interoperabilidad entre las distintas plantas del grupo Techint, como TenarisSiderca y TenarisDalmine. Esta homogeneidad tecnológica facilita el intercambio oportuno de información entre instalaciones, tanto para la mejora de procesos como para la prevención de problemas productivos en el conjunto de la organización.

Finalmente, la solución contribuye a fortalecer la escalabilidad y mantenibilidad de los sistemas informáticos de TenarisTamsa, alineándose con su modelo de excelencia operativa e innovación continua. Al reducir la dependencia de procesos manuales y proporcionar herramientas de visualización, proyección y análisis de información crítica, la organización estará mejor posicionada para atender el crecimiento constante de sus requerimientos operativos y para sostener su competitividad como proveedor de clase mundial.